\documentclass{article}
\usepackage[utf8]{inputenc}
\usepackage[T1]{fontenc}
\usepackage{geometry}
\usepackage{graphicx}
\usepackage{float}
\usepackage{hyperref}
\usepackage[english,indonesian]{babel}

\title{Optimasi Pemrosesan Template LaTeX Menggunakan Pendekatan Berbasis AI}
\author{Academic Support Agent}
\date{\today}

\begin{document}
\maketitle

\begin{abstract}

\end{abstract}

\chapter{PENDAHULUAN}


\chapter{TINJAUAN PUSTAKA}


% Original content (commented out):
% %%%%%%%%%%%%%%%%%%%%%%%%%%%%%%%%%%%%%%%%%%%%%%%%%%%%%%%%%%%%%%%%%%%%%%%
% Instructions: Comments in the file below indicate where the
% following steps have to be performed.
% Step 1: Enter abstract title.
% Step 2: Enter author information.
% Step 3: Enter key words.
% Step 4: Enter main text of abstract.
% Step 5: Enter references, e.g. using a simple list.
%%%%%%%%%%%%%%%%%%%%%%%%%%%%%%%%%%%%%%%%%%%%%%%%%%%%%%%%%%%%%%%%%%%%%%%
\documentclass[11pt,a4paper]{article}
\usepackage{times}
\usepackage{tgtermes}
\usepackage[T1]{fontenc}
% sur un systme mac os utiliser la ligne suivante
%\usepackage[applemac]{inputenc}
% sur un systme windows utiliser la ligne suivante
%\usepackage[latin1]{inputenc}
% sur un systme linux en utf8 utiliser al ligne suivante
\usepackage[utf8]{inputenc}
\usepackage[francais]{babel}
\usepackage{fancyhdr}
\usepackage{setspace}
\setstretch{1,15}

\usepackage{titlesec}
\titleformat{\section}
  {\normalfont\fontsize{16}{20}\bfseries}{\thesection}{1em}{}
\titleformat{\subsection}
  {\normalfont\fontsize{16}{20}\bfseries}{\thesubsection}{1em}{}
\titlespacing*{\section}
{0pt}{2.ex plus 1ex minus .2ex}{.0ex}
\titlespacing*{\subsection}
{0pt}{0.ex}{.0ex}

\setlength{\parindent}{0pt}
\setlength{\parskip}{5pt plus 2pt minus 1 pt}
\topmargin  -12mm
\evensidemargin 5mm
\oddsidemargin  0mm
\textwidth  158mm
\textheight 245mm
\headheight 14pt
\headsep 1.2cm

\pagestyle{fancy}
\rfoot{}
\chead{}
\cfoot{}
\lhead{
 \textit{23\textsuperscript{\`eme} Congr\`es Fran\c{c}ais de M\'ecanique}}              
 \rhead{
  \textit{Lille, 28 au 1\textsuperscript{er} Septembre 2017}}


\begin{document}
%\pagestyle{empty}



\begin{center}
\begin{spacing}{2.05}
{\fontsize{20}{20}
\bf
% Step 1: Enter abstract title here.
Titre\\ 
Times New Roman (20 gras)
}
\end{spacing}
\end{center}
\vspace{-1.25cm}
\begin{center}
{\fontsize{14}{20}
\bf
A. AUTEUR\textsuperscript{a}, A. AUTEUR\textsuperscript{b}, ... (14 gras)\\
\bigskip
%\vspace{0.75cm}
}
{\fontsize{12}{20}
a. Affiliation + email (12)\\
b. Affiliation + email (12)\\
...\\
}
\end{center}

\vspace{10pt}

{\fontsize{16}{20}
\bf
R\'esum\'e : (16 gras)
}
\medskip

\textit{Le Congr\`es Fran\c{c}ais de M\'ecanique a pour but principal de favoriser les 
\'echanges entre les acteurs de la communaut\'e qu'ils soient du monde acad\'emique 
ou du monde industriel, de pr\'esenter les challenges technologiques, d'identifier 
les probl\'ematiques scientifiques associ\'ees et de proposer des voies de progr\`es 
pour r\'epondre aux enjeux soci\'etaux. (11 Italique) Interligne (1.15)}

\vspace{20pt}

{\fontsize{16}{20}
\bf
Abstract : (16 gras)
}
\bigskip

\textit{The development of economical and safer societies requires to improve the 
means of transport and of energy production. Rotating machines hold a central place. 
(11 Italique) Interligne (1.15)}

\vspace{28pt}

{\fontsize{14}{20}
\bf
Mots clefs : Mots clefs (14 gras)
}
%\bigskip

\section{Introduction (16 gras)}
\medskip
Le Congr\`es Fran\c{c}ais de M\'ecanique a pour but principal de favoriser les 
\'echanges entre les acteurs de la communaut\'e qu'ils soient du monde acad\'emique 
ou du monde industriel... (11) Interligne (1.15)

\section{Paragraphe (16 gras)}

\subsection{Sous Paragraphe (16 gras)}
\medskip
Le Congr\`es Fran\c{c}ais de M\'ecanique a pour but principal de favoriser les 
\'echanges entre les acteurs de la communaut\'e qu'ils soient du monde acad\'emique 
ou du monde industriel... (11) Interligne (1.15)

%\section*{R\'ef\'erences (16 gras)}

\begin{thebibliography}{9}
\bigskip
  \bibitem{smith09}
	C. Smith, J.C. Green, 
	Modeled behavior of hydrodynamic bearings in thermal engine, 
	Mechanics \& Industry,
	10 (2009) 55--72.
  \bibitem{Ozisik73}
  	M.N. Ozisik,
	Radiative transfer and interactions with conduction and convection, 
	John Wiley and Sons, 
	New York, 
	1973.
  \bibitem{Sebastian92}
  	P. Sebastian,
	Analyse de réseaux de transferts : application à l'expertise de séchoirs.
	Thèse,
	Université de Bordeaux I,
	1992.
  \bibitem{Bougeard95}
  	D. Bougeard, J.-P. Vermeulen, B. Baudoin, 
	Spatial resolution enhancement of an IR system by image restoration techniques,
	in: D. Balageas, G. Busse, G.M. Carlomagno (ed.), 
	Proceedings of Quantitative Infrared Thermography QIRT 94 (Eurotherm Seminar 42),
	EDP Sciences,
	Paris, France, 
	1995, pp. 3--6.

\end{thebibliography}


%\begin{tabular}{lp{145mm}}
%% Step 5: Enter references here, e.g. using a simple list.
%\hspace{-2mm}{[1]}&
%N. Mo\"es, J. Dolbow and T. Belytschko.
%A finite element method for crack growth without remeshing. 
%{\em Int. J. Numer. Meth. Engng.}, Vol. {\bf 46}, 135--150, 1999.
%\\[1mm]
%
%\hspace{-2mm}{[2]}&
%O.C. Zienkiewicz and R.C. Taylor. 
%{\em The finite element method}, 
%4th. Edition, Vol. {\bf I}, McGraw Hill, 1989., Vol. {\bf II}, 1991.
%\\[1mm]
%\end{tabular}

\end{document}



\end{document}